\section*{الفصل \ref{ch:g2:born-grow-die} --- يولد وينمو ويموت}

\begin{solution}{exo:g2:born-grow-die:1}
\omterm{def:g2:born-grow-die:cycle}{الولادة}، والصغر، والبلوغ، \omterm{def:g2:born-grow-die:cycle}{والشيخوخة} (وتنتهي بالموت).
\end{solution}

\begin{solution}{exo:g2:born-grow-die:2}
في \omterm{def:g2:born-grow-die:cycle}{مرحلة} البلوغ.
\end{solution}

\begin{solution}{exo:g2:born-grow-die:3}
من الأكبر سنًّا إلى الأصغر: الجدة، ثم الأم، ثم التلميذ، ثم الرضيع.
\end{solution}

\begin{solution}{exo:g2:born-grow-die:4}
تنبت بذرة كرز (\omterm{def:g2:born-grow-die:cycle}{الولادة})؛ وتطول الشجرة الصغيرة سنوات؛ ثم تصنع
الشجرة البالغة الزهر والكرز --- ومن ثم بذورًا جديدة؛ ثم تشيخ
وتموت.
\end{solution}

\begin{solution}{exo:g2:born-grow-die:5}
مثلًا: الرضيع تنبت له أسنان، ثم تسقط، ثم تنبت أقوى منها؛ وزغب
الكتكوت الأصفر يصير ريشًا حقيقيًّا؛ \omterm{def:g1:plants-around-us:parts}{وساق} الشجرة الصغيرة الطرية
الخضراء تصير خشبًا صلبًا.
\end{solution}

\begin{solution}{exo:g2:born-grow-die:6}
من الأقصر إلى الأطول: الفراشة (أسابيع قليلة)، ثم القطة (من عشر
سنوات إلى عشرين)، ثم شجرة البلوط (مئات السنين).
\end{solution}

\begin{solution}{exo:g2:born-grow-die:7}
لأن للبالغين صغارًا قبل أن ينتهي الطريق، والصغار يبدؤون الطريق من
جديد --- وذلك هو السهم البنفسجي العائد. الحياة الواحدة طريق؛
والحياة عبر الأجيال تدور كالعجلة.
\end{solution}

\begin{solution}{exo:g2:born-grow-die:8}
الجواب مفتوح. والصور تظهر عادة الرضيع والطفل \omterm{def:g2:born-grow-die:cycle}{والبالغ} --- وقد تكون
\omterm{def:g2:born-grow-die:cycle}{الشيخوخة} هي \omterm{def:g2:born-grow-die:cycle}{المرحلة} التي لم تأتِ بعد.
\end{solution}

\begin{solution}{exo:g2:born-grow-die:9}
في تفاح الشجرة العجوز بذور، وكل بذرة تُزرع شجرة تفاح صغيرة: أي
الجيل التالي. وحين تموت الشجرة العجوز يكمل صغارها الدورة، فتحتفظ
الحديقة بأشجار تفاحها.
\end{solution}

\begin{solution}{exo:g2:born-grow-die:10}
لا --- بل لأن الحصاة \omterm{def:g1:living-or-not:nonliving}{لم تكن حية قط}. فالكائنات الحية وحدها تولد
وتنمو وتموت؛ والحصاة لا \omterm{def:g2:born-grow-die:cycle}{دورة حياة} لها البتة، فألّا تموت أبدًا ليس
قوةً، بل هو ألّا تكون حية.
\end{solution}
